\documentclass[12pt]{article}
\usepackage[utf8]{inputenc}
\usepackage[T1]{fontenc}
\usepackage{hyperref}
\title{Warnman Language Resources for K-12 Curriculum Development}
\author{Compiled by Claude Sonnet 5 (claude-sonnet-5) for Andr\'as Kornai}
\date{}
\begin{document}
\maketitle

\section{Introduction}
Warnman (Wanman; ISO 639-3 wbt; Glottocode wanm1242) is a Western Desert (Wati, Pama-Nyungan) language of the Jigalong area of Western Australia. It is spoken by the Warnman people, a subgroup of the Martu. Wikipedia gives 3 speakers (2016 census). A 2003 report mentions fewer than 100 speakers in the Pilbara, but that probably refers to a broader group. AIATSIS notes close relationships with Yulparija and Kukatja.

\section{Primary Educational Resources}
No school programme or primer was found.

\section{Linguistic and Reference Materials}
\begin{itemize}
\item Warnman language, Wikipedia. It cites Albert Burgman's \emph{Warnman dictionary} (Wangka Maya, 2010): \url{https://en.wikipedia.org/wiki/Warnman_language}
\item AIATSIS Austlang entry A62: \url{https://collection.aiatsis.gov.au/austlang/language/A62}
\item Mailing-list post (2006) on a Warnman--English wordlist and sketch grammar produced by Wangka Maya Pilbara Aboriginal Language Centre: \url{https://listserv.linguistlist.org/pipermail/ilat/2006-March/002255.html}
\item Vocabularies list, ANU: \url{https://www0.anu.edu.au/linguistics/nash/aust/dsd/vocabs.html}
\end{itemize}

\section{Religious and Community Resources}
Wangka Maya Pilbara Aboriginal Language Centre is named as the publisher of the wordlist and dictionary. I found no current website for it in this search, and the 2006 contact details may be outdated, so none are given here.

\section{Digital and Audio Resources}
The Endangered Languages catalogue was returned, but I did not confirm any Warnman entry: \url{https://endangeredlanguages.com/node?page=981}. Two Victorian Collections searches also appeared: \url{https://victoriancollections.net.au/collections?q=warnman}.

\section{Curriculum Development Context}
With very few speakers, any curriculum would be revitalisation-oriented and should be led by the Martu community and Wangka Maya, building on the 2003 wordlist and sketch grammar and the 2010 dictionary. The language is not known to be taught in schools.
\end{document}